%----------------------------------------------------------------------------------------
%	PACKAGES AND OTHER DOCUMENT CONFIGURATIONS
%----------------------------------------------------------------------------------------

\documentclass[12pt,stdletter,dateno,sigleft]{newlfm} % Extra options: 'sigleft' for a left-aligned signature, 'stdletternofrom' to remove the from address, 'letterpaper' for US letter paper - consult the newlfm class manual for more options

%\usepackage{charter} % Use the Charter font for the document text
\usepackage{graphics}

\usepackage{standalone}

%% with this patch, the newlfm should work (BEGIN)
\usepackage{etoolbox}
\makeatletter
\patchcmd{\@zfancyhead}{\fancy@reset}{\f@nch@reset}{}{}
\patchcmd{\@set@em@up}{\f@ncyolh}{\f@nch@olh}{}{}
\patchcmd{\@set@em@up}{\f@ncyolh}{\f@nch@olh}{}{}
\patchcmd{\@set@em@up}{\f@ncyorh}{\f@nch@orh}{}{}
\makeatother
%% with this patch, the newlfm should work (END)


\newsavebox{\Luiuc}\sbox{\Luiuc}{\parbox[b]{1in}{\vspace{0.35in}
\includegraphics[width=1.2\linewidth]{/Users/hectorbahamonde/Job_Market/logo.png}}} % Company/institution logo at the top left of the page
\makeletterhead{Uiuc}{\Lheader{\usebox{\Luiuc}}}

\newlfmP{sigsize=50pt} % Slightly decrease the height of the signature field

\lthUiuc % Print the company/institution logo

% email input begin
\newread\fid
\newcommand{\readfile}[1]% #1 = filename
{\bgroup
  \endlinechar=-1
  \openin\fid=#1
  \read\fid to\filetext
  \loop\ifx\empty\filetext\relax% skip over comments
    \read\fid to\filetext
  \repeat
  \closein\fid
  \global\let\filetext=\filetext
\egroup}
\readfile{/Users/hectorbahamonde/Bibliografia_PoliSci/email.txt}
% email input end

\usepackage{color} % for colors
\usepackage{hyperref} % for links

\hypersetup{
  colorlinks = true,
  urlcolor = blue,
  pdfpagemode = UseNone
}

%----------------------------------------------------------------------------------------
%	YOUR NAME AND CONTACT INFORMATION
%----------------------------------------------------------------------------------------

\namefrom{\vspace{-3cm}H\'ector Bahamonde, PhD} % Name

\addrfrom{
{\vspace{-2cm}}\\ % Date
[12pt]
{\bf H\'ector Bahamonde, PhD}\\
{\color{black}{\bf Profesor Asistente}}\\
\emph{{\color{black}Instituto de Ciencias Sociales}}\\
{\color{black}Universidad de O$'$Higgins}\\
Rancagua, Chile
}




%----------------------------------------------------------------------------------------
%	ADDRESSEE AND GREETING/CLOSING
%----------------------------------------------------------------------------------------

\greetto{\vspace{-1.9cm}Estimadas/os Integrantes de la Comisi\'on de Concursos Acad\'emicos,} % Greeting text
\closeline{\vspace{-1cm}\includegraphics[width=3cm]{/Users/hectorbahamonde/Administracion/signature.pdf} \hspace{-4cm}} % Closing text




%----------------------------------------------------------------------------------------

\begin{document}
\begin{newlfm}

%----------------------------------------------------------------------------------------
%	LETTER CONTENT
%----------------------------------------------------------------------------------------

Junto con saludarlos/as, me dirijo a ustedes para postular al concurso acad\'emico en temas del \'area de gobierno, gesti\'on y pol\'iticas p\'ublicas del Instituto de Asuntos P\'ublicos de la Universidad de Chile. 



Me motiva much\'isimo poder pertenecer al INAP ya que tiene uno de los mejores cuerpos acad\'emicos y estudiantes en el \'area de Administraci\'on P\'ublica del pa\'is. 

Aunque mi pregrado y doctorado son en ciencia pol\'itica, mis temas de investigaci\'on y docencia abordan aspectos inherentemente vinculados a temas del \'area de Gobierno y Administraci\'on P\'ublica. Esta es una de las razones por las que logr\'e asegurar una posici\'on a tiempo completo como profesor asistente (``\emph{tenure-track}'') en Administracion P\'ublica en la UOH. Siendo \emph{post-doctoral fellow} en el \emph{Center for Inter-American  Policy and Research} (Tulane University) fui convocado a ser el {\bf primer acad\'emico \emph{full-time} que asumi\'o el rol de Jefe de Carrera de Administraci\'on P\'ublica UOH}. Mientras cumpl\'ia este rol (cerca de un a\~no) tuve la oportunidad de construir universidad p\'ublica, abrir nuevas fronteras, asegurar una educaci\'on p\'ublica de calidad a las y los estudiantes de la regi\'on, y de colaborar en la descentralizaci\'on el pa\'is. Y aunque ya no cumplo ese rol debido a razones de reestructuramiento organizacional (ahora las/los jefes de carreras son medias jornadas sin doctorado), mantengo estrecha colaboraci\'on con la actual Jefa de Carrera y sigo ense\~nando a las/los futuras/os administradoras/es p\'ublicas/os de la Regi\'on de O'Higgins. {\bf Ellas/os siguen siendo mi principal \emph{constituency}}. Es importante notar que gran parte de mi \emph{teaching} est\'a orientado a administradores/as p\'ublicos/as. Por ejemplo, en \href{https://github.com/hbahamonde/Ciencia_Politica_I/raw/master/Bahamonde_Ciencia_Politica_I.pdf}{\emph{Ciencia Pol\'itica II}} ense\~no formaci\'on del aparato p\'ublico-administrativo y en \href{https://github.com/hbahamonde/Ciencia_Politica_I/raw/master/Bahamonde_Ciencia_Politica_I.pdf}{\emph{M\'etodos de Investigaci\'on}} ense\~no an\'alisis de datos para administradores/as p\'ublicos/as.


Mi tema principal de investigaci\'on es la {\bf formaci\'on de capacidades administrativas del estado Chileno}. En mi tesis doctoral---as\'i como en los \href{https://github.com/hbahamonde/Earthquake_Paper/raw/master/Bahamonde_Earthquake_Paper.pdf}{\emph{papers}} que pueden ver en mi CV---estudio los mecanismos por los cuales los estados en Latinoam\'erica y en Chile construyen (o no) capacidades administrativas para hacer respetar normativas de construcci\'on (casas, puentes, etc.). Un pa\'is aplanado por un terremoto no es otra cosa que una cordial invitaci\'on al estado natural de Hobbes. En tanto pol\'itica p\'ublica, yo explico que aunque todos los estados \emph{quieren} tener las capacidades administrativas para hacer cumplir esta pol\'itica p\'ublica (y evitar una cat\'astrofe), no todos \emph{pueden}. Parte de la explicaci\'on es que los estados que presentan disputas sectoriales dentro de sus econom\'ias logran movilizar los incentivos necesarios para crear un aparato fiscal fuerte. Levantando el marco te\'orico de la ``sociolog\'ia fiscal'', yo explico c\'omo los impuestos crearon al estado y c\'omo los estados crearon los impuestos. En consecuencia, parte de mi contribuci\'on a la administraci\'on p\'ublica como disciplina consiste en formular condiciones estructurales que ayudaron a los estados Latinoamericanos a desarrollar sus capacidades administrativas y de pol\'itica p\'ublica durante la primera mitad del siglo XIX. 


Mi segundo tema de investigaci\'on examina cu\'ales son las fallas administrativas y de pol\'itica p\'ublica del Estado que contribuyen a la aparici\'on del {\bf clientelismo}. Quisiera nombrar dos ejemplos. En mi \emph{paper} ``\href{https://journals.sub.uni-hamburg.de/giga/jpla/article/view/1121/1128}{Aiming Right at You}'' (WoS) me focalizo en c\'omo la pol\'itica p\'ublica del urbanismo podr\'ia potenciar el clientelismo en Brasil. En particular, el argumento explica c\'omo situaciones de heterogeneidad en niveles de desigualdad econ\'omica aumenta la visibilidad de hogares dispuestos a vender su voto. El segundo ejemplo de c\'omo mi investigaci\'on en clientelismo toca aspectos fundamentales relacionados a las capacidades administrativas del estado y la pol\'itica p\'ublica, es mi \emph{paper} ``\href{https://doi.org/10.1057/s41269-020-00174-4}{Still for Sale}'' (WoS). En este trabajo estudio el caso de c\'omo los habitantes de un estado administrativamente fuerte y cap\'az como Estados Unidos son susceptibles de incurrir en pr\'acticas clientelares. El argumento que desarrollo propone ciertas condiciones individuales bajo las cuales los ciudadanos estar\'ian dispuestos a vender su voto. Tal como lo explico en esta \href{https://eltipografo.cl/2019/07/desigualdad-y-democracia-en-chile}{columna de opini\'on}, las implicancias en t\'erminos de pol\'iticas p\'ublicas son muy grandes, sobre todo en t\'erminos de la pol\'itica p\'ublica del financiamiento de campa\~nas.



Otro tema de investigaci\'on que investigo, relativo a la administraci\'on y a las pol\'iticas p\'ublicas del pa\'is, es la {\bf descentralizaci\'on}. Por ejemplo, en mi \emph{paper} ``\href{https://doi.org/10.1111/rsp3.12337}{Employment Effects of Covid‐19 across Chilean Regions}'' (WoS) estudiamos c\'omo la pandemia del COVID-19 golpe\'o de distintas maneras al pa\'is: usando t\'ecnicas de \emph{forecasting}, encontramos que las regiones mineras tienden a ser m\'as resistentes---en t\'erminos de p\'erdida de empleos---al ser comparadas con regiones agr\'icolas y forestales. La tensi\'on entre descentralizaci\'on y lo estatal me causan mucho inter\'es: en este minuto me encuentro desarrollando un experimento para estudiar los niveles sub-nacionales de {\bf gobernanza}. Mi pregunta es {\bf \emph{¿penetran las pol\'iticas p\'ublicas del Estado todas las geograf\'ias por igual?}}


Para cerrar esta carta, me gustar\'ia volver a expresar mi inter\'es en este concurso. Gran parte de los temas que desarrollo est\'an vinculados a la administraci\'on p\'ublica, y en ese sentido, insertarme en un ambiente acad\'emico dotado de esta \emph{expertise} contribuir\'a a mejorar mi investigaci\'on en estas \'areas. Me encantar\'ia poder colaborar con muchos de los/las actuales investigadores/as, como por ejemplo, Letelier, Pliscoff y Figueroa, entre otros/as. Tambi\'en me interesa colaborar con estudiantes de postgrado. Como lo se\~nala mi CV, soy un fiel creyente en la colaboraci\'on interdisciplinar. Adem\'as, me sentir\'ia muy c\'omodo ense\~nando cursos de formaci\'on del estado, administraci\'on p\'ublica y fen\'omenos pol\'iticos, an\'alisis emp\'irico de pol\'itica p\'ublica, epistemolog\'ia, metodolog\'ia (cuali/cuanti), entre otros. 

Junto con despedirme, invito a las/los miembros del \emph{Comit\'e} a considerar mi postulaci\'on.\unskip



%----------------------------------------------------------------------------------------

\end{newlfm}
\end{document}
